\begin{abstract}
We present simulation-based inference of thermal Sunyaev--Zel'dovich (tSZ) cluster counts in which every
training simulation runs the cluster finder on simulated multi-frequency maps, to our knowledge the first
of its kind. Our simulator paints haloes into six Simons Observatory (SO) frequency maps with the cosmic
microwave background (CMB) and goal-depth noise, and runs a matched-filter finder. On one GPU, a survey of
40.7 per cent of the sky takes 1.6\,s and a training set of 10\,000 surveys 4.6\,h. At this sky fraction, the
posterior is within 0.11 widths (posterior standard deviations) of the truth, and its 90 per cent
intervals contain the truth for 0.88--0.91 of 500 prior skies. On the same skies, we compare it with
catalogue models calibrated on map-level mocks, which have to carry the finder's confusion noise (4--8 per
cent), search over positions and sizes, neighbours and upward error tail. A model calibrated at the
true parameters is within 0.2 widths at that point in $\sigma_8$, the cold dark matter density and the
tSZ--mass log-amplitude, but too narrow or off-centre across the prior. Calibrated at 25 points across
the prior, it has close to correct coverage for these three but leaves the slope of the relation
0.6--0.7 widths high, which we attribute to the tail. As a reference, the uncalibrated model, with a noise
table for the instrument and the CMB only, is biased by up to 8.3 widths on our skies, showing the size of
the finder's effects at goal depth.
\end{abstract}
